\begin{equation}
 \sum_{\sigma=-s}^{+s}|\psi(\sigma)|^2
 =\sum_{\lambda,\mu,\ldots=1}^{2}|\psi^{\lambda\mu\ldots}|^2
 \tag{57.1}
\end{equation}
(This sum is a scalar, as it must be, since it gives the probability of finding
the particle at the specified point in space.) In the sum on the right-hand
side, a component having $(s+\sigma)$ indices equal to 1 occurs
\[
 \frac{(2s)!}{(s+\sigma)!(s-\sigma)!}
\]
times. It follows that the correspondence between the functions $\psi(\sigma)$
and the spinor components is given by
\begin{equation}
 \psi(\sigma)=\left[\frac{(2s)!}{(s+\sigma)!(s-\sigma)!}\right]^{1/2}
 \psi^{\overbrace{11\ldots1}^{s+\sigma}
       \overbrace{22\ldots2}^{s-\sigma}}.
 \tag{57.2}
\end{equation}
